\begin{abstract}
Cardiovascular disease is the leading cause of death worldwide, and most machine-learning work on cardiac screening uses a single data source. MultiCardioFusion enriches the 13 clinical variables of the UCI Cleveland Heart Disease dataset with 23 ECG-proxy features derived from the same patients' resting-ECG and exercise-test measurements, benchmarks a parallel PTB-XL ECG branch with 568 tsfresh descriptors, and explains predictions with SHAP. Because Cleveland and PTB-XL contain different patients, the work is framed as ECG-proxy-enhanced clinical prediction, not patient-level multimodal fusion. An earlier version reported AUC 0.9535 and a 7.04-point accuracy gain from fusion on a single 61-patient hold-out. Here every claim is re-tested under a controlled protocol: all preprocessing inside training folds, 5$\times$10 repeated cross-validation with identical splits and tuning budget, corrected confidence intervals and Holm-adjusted tests, over \nVariants{} configurations. These include component removal, negative controls, a permutation test, eight imbalance strategies, leakage and hold-out sensitivity analyses, 16 robustness scenarios, subgroups, external validation on three hospitals, SHAP stability, cross-attention variants and a PTB-XL lead and feature ablation. The full system reaches AUC \fullAUC{} \fullAUCci{}, but the clinical variables alone perform equally ($\Delta$AUC \dClinAUC{} in favour of clinical-only). The proxy block performs no better than added noise. The reported gain arose from subtracting a cross-validation accuracy from a hold-out accuracy of a different model, on a split that lies at the 95th percentile of \hoN{} random splits. SMOTE and length-1 cross-attention add nothing, oversampling before cross-validation inflates AUC, and external AUC ranges from \extVaFu{} to \extHunFu{}. SHAP rankings of \texttt{cp}, \texttt{ca} and \texttt{thal} are stable, while credit given to engineered features is redistributed from their parent variables. Because the dominant predictors come from fluoroscopy and thallium scintigraphy, the model is a post-test classifier and is not validated for clinical use. The study provides a leakage-aware, fully ablated baseline and shows how evaluation choices can create apparent fusion gains on small clinical datasets.
\end{abstract}

\begin{IEEEkeywords}
heart disease prediction, ablation study, ECG-proxy features, repeated cross-validation, data leakage, SHAP stability, external validation, calibration, tsfresh, PTB-XL.
\end{IEEEkeywords}
